\documentclass[10.5pt]{config}
% 本实验报告采用 署名-相同方式共享 4.0 国际许可协议 (CC BY-SA 4.0)

% 进一步更改，如实验名称后面横线的长度等，需在config.cls文件中修改

% ===================== 使用说明 =====================
% 1. 编译引擎必须为 XeLaTeX：latexmk -xelatex main.tex（执行两遍）。
% 2. 第一次报告：打印《00 1 实验报告封面.doc》填写后装订在本 PDF 之前；
%    第 2 次起不再需要封面，直接用 \makeheader 生成的页眉。
% 3. 打印稿首页左上角手写席卡序号（实验台席卡上姓名右侧数字）。
% 4. 实验名称、实验类型逐字取自讲义/当学期通知，不得简写。
% 5. 模板中的指导性注释不要删除；内容写在注释下方。
% ====================================================

\major{}
\name{} % 你的姓名
\partner{} % 组员姓名
% \title{}
\stuid{} % 你的学号
%\partnerstuid{} % 组员学号
\date{} % 实验日期，以第一次为例
\lab{} % 实验地点
\course{电子电路设计实验I} % 课程名称
\instructor{} % 指导老师
\grades{} % 成绩，应该不填，此处按老师提供的模板进行保留
\expname{集成运算放大器应用电路研究（II）} % 实验名字，以第一次为例
\exptype{设计类实验}
\begin{document}

% \maketitlepage % 封面，原latex模板中可选择进行配置
% \maketoc % 目录，说明同上
\makeheader % 首页头部


\section{实验目的}
% 3–5 条，动词开头；取自讲义对应实验的“实验目的”，不扩写。

\section{实验任务与要求}
% 抄全讲义任务，并量化指标：输入幅度/频率、增益、误差要求、待验证的定理等。

\section{实验原理}
% 必要公式 + 电路图/框图；给出推导关键步与测量原理，不堆砌教材内容。

\section{实验方案设计与参数计算}
\subsection{实验总体方案设计}
% 总体框图、信号流向、方案选择理由。
\subsection{各功能电路设计与计算}
% 元件参数选择与理论计算（网孔/节点法），列理论值表。
\subsection{完整的实验电路}
% 完整电路原理图 + 实物接线图；标注测试点、电源与地。

\section{实验仪器设备}
% 仪器名称、型号、关键参数；软件写版本（如 OrCAD 26.1 / PSpice AD）。

\section{实验步骤、调试过程、实验数据记录}
% 按时间顺序写；故障现象 → 判断 → 处理；原始数据表（单位、小数位统一）；
% 仿真设置与波形截图；记录本照片为数据来源时，转录后需与用户核对。

\section{数据分析与讨论}
% 评分权重最高（25%）。至少包含：
% 1) 理论值/实测值/绝对误差/相对误差表；
% 2) 关键关系图（U-I、幅频、波形对比），含拟合线或理论线；
% 3) 误差来源逐条解释（仪表内阻、元件容差、接触电阻、电源漂移、单次测量、仿真理想化）；
% 4) 异常现象的解释与验证；
% 5) 设计参数的修正建议（如有）。

\section{结论}
% 编号列出，逐条对应“实验任务与要求”；带定量结果；不引入新数据/新图。

\section{心得与体会}
% 写具体收获：调试方法、故障排查、软件使用、理论联系实际。

\section{思考题}
% 逐题作答，写推理与公式，不跳步。

\end{document}
